\documentclass[a4paper,11pt]{jsarticle}
\usepackage[dvipdfmx]{graphicx}
\usepackage{geometry}
\usepackage{amsmath}
\usepackage{booktabs}
\usepackage{url}
\geometry{margin=20mm}

\title{ACC 時間項 \(p_t(t)\) の扱いに関する補足ノート\\
AW 除外時に \(\int_{\mathrm{AW}} \mathrm{tshape}\,dt = 0\) となる理由}
\author{p2MEG 解析メモ}
\date{2026年3月8日}

\begin{document}
\maketitle

\section{このノートの目的}
本ノートの目的は、現在の ACC PDF で時間項がどのように扱われているか、
特に
\begin{quote}
「今回は 4D AW 内事象を time template からも除外したため、
保存された tshape の AW 内積分は 0 になった」
\end{quote}
という文の意味を、図を使って具体的に説明することである。

\section{まず結論}
今のコードでは、ACC の時間テンプレート \texttt{tshape} は
\[
E_e \in \mathrm{AW},\quad
\theta_{eg}\in \mathrm{AW},\quad
E_\gamma \in [20,30]\ \text{or}\ [30,80]\ \mathrm{MeV}
\]
の事象から作る。
ただし、さらに
\[
(E_e,E_\gamma,\theta_{eg},t)\in \mathrm{AW}
\]
の事象は埋めない。

今回の \(E_\gamma\) template bin は \([20,30]\) と \([30,80]\) MeV なので、
これはどちらも最終解析窓の \(E_\gamma\in[20,80]\) MeV の中に入っている。
したがって、time template に対して AW 除外をかけると、
実質的には
\[
t\in[-50,50]\ \mathrm{ns}
\]
の領域に事象が全く入らなくなる。

その結果、保存された \texttt{tshape} は TSB にだけ中身を持ち、
AW では 0 になる。
ゆえに
\[
\int_{-50}^{50} \mathrm{tshape}(t)\,dt = 0
\]
となる。

\section{コード上で何をしているか}
\subsection{time template を埋める条件}
\texttt{src/MakeACCGridPdf.cc} では、time template は概念的に次の条件で埋めている。
\begin{enumerate}
  \item \(t \in [-500,500]\) ns
  \item \(E_e\) は AW 内
  \item \(\theta_{eg}\) は AW 内
  \item \(E_\gamma\) は time template 用カテゴリ \([20,30]\) または \([30,80]\) MeV
  \item ただし 4D AW 内 \((E_e,E_\gamma,\theta_{eg},t)\in\mathrm{AW}\) は除外
\end{enumerate}

最後の条件のため、今の template は
\[
t \in [-500,-50)\cup(50,500]
\]
にしか実質的に事象を持たない。

\subsection{なぜ AW 積分が 0 になるのか}
\texttt{tshape} は parametric fit 関数ではなく、単なる 1D ヒストグラム密度である。
したがって、AW に一度も fill されていないなら、
AW 内の各 bin content は 0 である。

従って、
\[
\int_{-50}^{50} \mathrm{tshape}(t)\,dt
\]
をヒストグラム密度として積分すると 0 になる。

ここで重要なのは、
\begin{quote}
``TSB にだけデータがあるので、そこから自動で AW に滑らかに延長してくれる''
\end{quote}
わけではない点である。
今の \texttt{MakeACCGridPdf} は time template に対しては
そのような関数外挿をしておらず、単に TSB だけ埋まったヒストを保存している。

\section{ACCGridPdf 側で何が起きるか}
\texttt{src/ACCGridPdf.cc} では、ロード時に
\[
\mathrm{AW\ norm}
=
\int_{\mathrm{AW}} \mathrm{tshape}(t)\,dt
\]
を計算し、この値で AW 内の時間 PDF を正規化しようとする。

しかし今回の \texttt{tshape} ではこの積分が 0 なので、
\begin{enumerate}
  \item Eg ごとの \texttt{tshape\_egbin*} を無効
  \item 全体 \texttt{tshape} も無効
  \item 最後に AW 内一様分布 \(1/(t_{\max}-t_{\min})\) へフォールバック
\end{enumerate}
となる。

したがって、今の ACC PDF の時間項は最終的に
\[
p_t(t)=\frac{1}{100\ \mathrm{ns}}
\qquad
(-50 \le t \le 50\ \mathrm{ns})
\]
である。

\section{なぜこうしたのか}
これは
\begin{quote}
``AW の事象を ACC モデルに使わない''
\end{quote}
という方針を徹底した結果である。

もし AW 内時間形状もヒストグラムとして持ちたいなら、次のどちらかが必要になる。
\begin{enumerate}
  \item AW 内事象を template に入れる
  \item AW は入れずに、TSB を関数で fit して AW に外挿する
\end{enumerate}

今の ACC PDF 生成器は 2 番目をまだ実装していない。
だから、AW を除外したままではヒストグラム template は AW で 0 になり、
結果として一様時間にフォールバックする。

\section{図で見る}
\subsection{ACC 予測の図}
まず、TSB から AW 内 ACC 数を予測した結果を図\ref{fig:pred1}に示す。

\begin{figure}[p]
  \centering
  \includegraphics[width=0.96\linewidth,page=1]{predict_acc_from_tsb_by_eg_step4f_run1to20_eg_doubleonly_allpatterns_500ns.pdf}
  \vspace{3mm}
  \includegraphics[width=0.96\linewidth,page=2]{predict_acc_from_tsb_by_eg_step4f_run1to20_eg_doubleonly_allpatterns_500ns.pdf}
  \caption{TSB から AW 内 ACC 数を外挿した図。}
  \label{fig:pred1}
\end{figure}

\subsection{時間 fit の図}
AW 除外後の時間分布 fit を図\ref{fig:fitall}に示す。
これらの図で使っている \(t\) ヒストグラムには、4D AW 内事象は入っていない。

\begin{figure}[p]
  \centering
  \includegraphics[width=0.96\linewidth,page=2]{fit_acc_time_by_eg_step4f_run1to20_eg_doubleonly_allpatterns_500ns.pdf}
  \vspace{2mm}
  \includegraphics[width=0.96\linewidth,page=3]{fit_acc_time_by_eg_step4f_run1to20_eg_doubleonly_allpatterns_500ns.pdf}
  \caption{\(E_\gamma=[0,10]\), \([10,20]\) MeV の時間 fit。}
  \label{fig:fitall}
\end{figure}

\begin{figure}[p]
  \centering
  \includegraphics[width=0.96\linewidth,page=4]{fit_acc_time_by_eg_step4f_run1to20_eg_doubleonly_allpatterns_500ns.pdf}
  \vspace{2mm}
  \includegraphics[width=0.96\linewidth,page=5]{fit_acc_time_by_eg_step4f_run1to20_eg_doubleonly_allpatterns_500ns.pdf}
  \caption{\(E_\gamma=[20,30]\), \([30,80]\) MeV の時間 fit。}
  \label{fig:fitall2}
\end{figure}

\subsection{ACC PDF 本体の図}
現在の \texttt{acc\_grid.root} を可視化したものを図\ref{fig:accgridmeta}と図\ref{fig:accgridshape}に示す。
meta page に
\begin{quote}
\texttt{t shape: uniform fallback in AW}
\end{quote}
と出ているのが、今回の議論の直接の証拠である。

\begin{figure}[p]
  \centering
  \includegraphics[width=0.96\linewidth,page=1]{acc_grid_hist_step4f_doubleonly_awexcluded.pdf}
  \caption{ACC PDF 可視化の meta page。\texttt{t shape: uniform fallback in AW} と表示されている。}
  \label{fig:accgridmeta}
\end{figure}

\begin{figure}[p]
  \centering
  \includegraphics[width=0.96\linewidth,page=2]{acc_grid_hist_step4f_doubleonly_awexcluded.pdf}
  \vspace{2mm}
  \includegraphics[width=0.96\linewidth,page=3]{acc_grid_hist_step4f_doubleonly_awexcluded.pdf}
  \caption{ACC 4D PDF の 1D/2D 可視化。ここで見えている形状は 4D の \(p_4(E_e,E_\gamma,\phi_e,\phi_\gamma)\) に由来する。時間項そのものは AW 内一様へフォールバックしている。}
  \label{fig:accgridshape}
\end{figure}

\section{実務的な理解}
今回の実装を一言で言えば、
\begin{quote}
ACC のエネルギー・角度形状は TSB だけで作った。
時間形状も AW を使わずに作ろうとしたが、ヒストグラム法では AW に中身が残らないので、
最終的には時間一様に戻った。
\end{quote}
ということである。

つまり、今の尤度 fit で使っている ACC PDF は
\[
p_{\mathrm{acc}}(x)
=
p_4(E_e,E_\gamma,\phi_e,\phi_\gamma)\times \mathrm{const.}
\]
であり、時間構造を AW 内で能動的に学習した PDF にはまだなっていない。

\section{次にもし改善するなら}
AW を見ずに時間項もデータ駆動で持ち込みたいなら、
次の方法が自然である。
\begin{enumerate}
  \item 今の \texttt{fit\_acc\_time\_by\_eg.C} のように TSB を関数 fit する
  \item その fitted function を \texttt{acc\_grid.root} に関数パラメータまたは AW 正規化済みヒストとして保存する
  \item \texttt{ACCGridPdf} 側でその関数を使って \(p_t(t|E_\gamma)\) を評価する
\end{enumerate}

そうすれば、
\begin{quote}
AW 内事象を使わず、かつ AW 内で 0 にならない時間項
\end{quote}
を作ることができる。

\end{document}
